% Format teze zasnovan je na paketu memoir
% http://tug.ctan.org/macros/latex/contrib/memoir/memman.pdf ili
% http://texdoc.net/texmf-dist/doc/latex/memoir/memman.pdf
% 
% Prilikom zadavanja klase memoir, navedenim opcijama se podešava 
% veličina slova (12pt) i jednostrano štampanje (oneside).
% Ove parametre možete menjati samo ako pravite nezvanične verzije
% mastera za privatnu upotrebu (na primer, u b5 varijanti ima smisla 
% smanjiti 
\documentclass[12pt,oneside]{memoir}

% Paket koji definiše sve specifičnosti mastera Matematičkog fakulteta
\usepackage{matfmaster}
%
% Podrazumevano pismo je ćirilica.
%   Ako koristite pdflatex, a ne xetex, sav latinički tekst na srpskom jeziku
%   treba biti okružen sa \lat{...} ili \begin{latinica}...\end{latinica}.
%
% Opicija [latinica]:
%   ako želite da pišete latiniciom, dodajte opciju "latinica" tj.
%   prethodni paket uključite pomoću: \usepackage[latinica]{matfmaster}.
%   Ako koristite pdflatex, a ne xetex, sav ćirilički tekst treba biti
%   okružen sa \cir{...} ili \begin{cirilica}...\end{cirilica}.
%
% Opcija [biblatex]:
%   ako želite da koristite reference na više jezika i umesto paketa
%   bibtex da koristite BibLaTeX/Biber, dodajte opciju "biblatex" tj.
%   prethodni paket uključite pomoću: \usepackage[biblatex]{matfmaster}
%
% Opcija [b5paper]:
%   ako želite da napravite verziju teze u manjem (b5) formatu, navedite
%   opciju "b5paper", tj. prethodni paket uključite pomoću: 
%   \usepackage[b5paper]{matfmaster}. Tada ima smisla razmisliti o promeni
%   veličine slova (izmenom opcije 12pt na 11pt u \documentclass{memoir}).
%
% Naravno, opcije je moguće kombinovati.
% Npr. \usepackage[b5paper,biblatex]{matfmaster}

% Paket koji obezbeđuje ispravni prikaz ćiriličkih italik slova kada
% se koristi pdflatex. Zakomentarisati ako na sistemu koji koristite ovaj
% paket nije dostupan ili ako ne radi ispravno.
\usepackage{cmsrb}

% Ostali paketi koji se koriste u dokumentu
\usepackage{listings} % listing programskog koda

% Vidljiva oznaka za planirane slike u radnoj verziji. Pri zameni je dovoljno
% zameniti poziv ove komande sa \includegraphics; natpis i oznaka ostaju isti.
\newcommand{\draftfigureplaceholder}[2]{%
	\centering
	\begingroup
	\setlength{\fboxsep}{10pt}%
	\fbox{%
		\parbox{0.88\textwidth}{%
			\centering\textbf{МЕСТО ЗА #1}\par
			\medskip
			\raggedright\small\textbf{Предвиђени садржај:} #2%
		}%
	}%
	\endgroup
}

% Datoteka sa literaturom u BibTex tj. BibLaTeX/Biber formatu
\bib{matf-master-pavle-vlajkovic}

% Ime kandidata na srpskom jeziku (u odabranom pismu)
\autor{Павле Влајковић}
% Naslov teze na srpskom jeziku (u odabranom pismu)
% TODO: Са ментором потврдити коначни наслов рада.
\naslov{Емпиријско поређење архитектуралних образаца софтверских система на примеру обраде мултимедијалног садржаја}
% Godina u kojoj je teza predana komisiji
\godina{2026}
% Ime i afilijacija mentora (u odabranom pismu)
% TODO: После званичне провере допунити титулу и афилијацију ментора.
\mentor{Владимир Филиповић}
% Ime i afilijacija prvog člana komisije (u odabranom pismu)
% TODO: После званичне провере допунити титулу и афилијацију члана комисије.
\komisijaA{Александар Картељ}
% Ime i afilijacija drugog člana komisije (u odabranom pismu)
% TODO: После званичне провере допунити титулу и афилијацију члана комисије.
\komisijaB{Сташа Вујчић Станковић}
% Ime i afilijacija trećeg člana komisije (opciono)
% \komisijaC{}
% Ime i afilijacija četvrtog člana komisije (opciono)
% \komisijaD{}
% Datum odbrane (obrisati ili iskomentarisati narednu liniju ako datum odbrane nije poznat)
% \datumodbrane{}

% Apstrakt na srpskom jeziku (u odabranom pismu)
% TODO: Написати апстракт тек после усвајања коначног обима и резултата.
\apstr{}

% Ključne reči na srpskom jeziku (u odabranom pismu)
% TODO: Одредити кључне речи после усвајања коначног обима рада.
\kljucnereci{}

\begin{document}
% ==============================================================================
% Uvodni deo teze
\frontmatter
% ==============================================================================
% Naslovna strana
\naslovna
% Strana sa podacima o mentoru i članovima komisije
\komisija
% Strana sa posvetom (u odabranom pismu)
% TODO: Потврдити да ли посвета остаје у коначној верзији.
\posveta{Мами, тати и деди}
% Strana sa podacima o disertaciji na srpskom jeziku
\apstrakt
% Sadržaj teze
\tableofcontents*

% ==============================================================================
% Glavni deo teze
\mainmatter
% ==============================================================================

% ------------------------------------------------------------------------------
\chapter{Увод}
\label{chp:uvod}
% ------------------------------------------------------------------------------
\begingroup
\setlength{\emergencystretch}{3em}

\section{Предмет и мотивација рада}

Софтверска архитектура одређује основну организацију система, односе међу
његовим деловима и оквир за касније пројектне одлуке. Она је основа за
разматрање атрибута квалитета и обухвата ране одлуке које усмеравају развој
система \cite{bass2012software}. Избор између монолита, модуларног монолита и
микросервисне архитектуре зато зависи од конкретног домена, оптерећења и
операционих услова.

Предмет рада је поређење три архитектуралне варијанте система за обраду видео
садржаја. Систем прихвата видео датотеку, издваја метаподатке, формира послове
за израду варијанти у задатим резолуцијама и обавља
транскодовање\footnote{Транскодовање је промена начина кодовања или параметара
	медијског записа.} помоћу алата {\lat FFmpeg}. Посматрани ток укључује веб
интерфејс, трајно чување података, асинхрону обраду, ред послова, спољни процес
и складиште датотека. Због тога омогућава испитивање архитектуралних разлика у
вишефазној обради, а не само у једноставним синхроним операцијама.

% Извор описа система: backend/media-monolith и backend/media-modular;
% функционални ток: benchmark/gatling/src/test/java/benchmark/LoadStressSimulation.java;
% обрада послова: одговарајуће датотеке RenditionJobHandler.java у обе варијанте.

Практичне препоруке, попут предлога да се нов систем под одређеним условима
најпре развија као монолит \cite{fowler2015monolithfirst}, не одговарају саме
по себи на питање како се три варијанте овог система понашају под истим
условима. Рад се стога бави системски специфичним недостатком доказа: за
посматрани ток обраде, технолошки скуп и експериментално окружење потребно је
контролисано квантитативно и квалитативно поређење, праћено евиденцијом промена
током миграције.

\section{Циљ и истраживачка питања}

Циљ рада је да се на истом функционалном примеру испита однос архитектуралне
организације, мерљивих перформанси и промена насталих миграцијом од монолита,
преко модуларног монолита, до микросервиса. Прва етапа миграције уводи јасне и
аутоматски проверљиве границе модула у заједничкој јединици распоређивања.
Друга етапа издваја пословне компоненте у сервисе који се могу независно
распоређивати.

Истраживање одговара на три питања:
\begin{enumerate}
	\item[ИП1.] Какве се разлике у пропусности, времену одзива, латенцији тока
		обраде, употреби процесора и употреби меморије јављају између три
		варијанте при истом оптерећењу и у дефинисаним ресурсним
		профилима\footnote{Ресурсни профил је унапред одређена комбинација
			укупног буџета процесора и меморије и дозвољене конкурентности
			апликационог слоја.}?
	\item[ИП2.] Које се разлике у организацији кода, сложености распоређивања и
		понашању при отказу уочавају применом истих критеријума, артефаката и
		проба на све три варијанте?
	\item[ИП3.] Које врсте документованих промена настају у свакој етапи
		миграције и какав је њихов описни обим, исказан бројем јединица по
		категорији и категоријама захваћених артефаката?
\end{enumerate}

Пропусност\footnote{Пропусност је број успешно завршених пословних операција у
	јединици времена; и операција и мерни интервал морају бити унапред одређени.}
и латенција тока обраде\footnote{Латенција тока обраде мери време од креирања
	посла до његовог коначног исхода и обухвата чекање у реду.} односе се на цео
пословни ток, док време одзива описује појединачни {\lat HTTP} захтев. Ове мере
се зато одређују и тумаче засебно. Њихове операционе дефиниције морају бити
једнаке за све варијанте \cite{jain1991art,kounev2020benchmarking}.

Препоруке за избор архитектуре нису засебно истраживачко питање. Оне се у
дискусији и закључку изводе условном синтезом налаза ИП1--ИП3, у складу са
приоритетима и ограничењима посматраног случаја.

\section{Обим и приступ}

Студија је контролисано поређење у коме је архитектурална варијанта фактор, а
монолит, модуларни монолит и микросервисна архитектура његова три нивоа.
Варијанте чувају исти домен, функционални обим, јавни програмски интерфејс,
логику транскодовања, тест податке, сценарио оптерећења и окружење мерења.
Намерно се разликују модулске и процесне границе, комуникација и топологија
распоређивања. Методологија мерења захтева да систем под тестом, фактор,
контролне променљиве, оптерећење и метрике буду унапред одређени
\cite{jain1991art,kounev2020benchmarking}.

Квантитативни део обухвата независно поновљена мерења под заједничким
ресурсним буџетом. Квалитативни део примењује симетричне поступке на
организацију кода, распоређивање и отказе. Анализа миграције користи
верзионисане промене и друге проверљиве артефакте; без посебних временских
записа из њих се не изводе тврдње о уложеном раду, трајању или трошку.

Обим је ограничен на транскодовање видео-записа, изабрани технолошки скуп,
синтетичке тест-податке, дефинисано оптерећење и локално окружење. Израда
слика, издвајање звука, продукционо оптерећење и оркестрирано окружење нису
предмет експеримента. Ова ограничења одређују и домет завршних препорука.

\section{Намеравани допринос и структура рада}

Намеравани допринос је упоредив скуп квантитативних података за три варијанте,
симетрично квалитативно поређење и документована двоетапна миграција. Њихова
синтеза треба да омогући условне препоруке за системе са сродним захтевима и
условима.

% TODO: Пре коначне предаје потврдити да су микросервисна варијанта, друга
% етапа миграције и комплетна мерна матрица завршене и архивиране; затим
% формулацију намераваног доприноса ускладити са стварно доказаним резултатима.

После увода, друго поглавље утврђује појмове и критеријуме класификације, а
треће описује метод истраживања. Поглавља четири, пет и шест приказују три
варијанте и две етапе миграције. Затим следе резултати и дискусија, однос према
сродним радовима и закључак.
\endgroup

% ------------------------------------------------------------------------------
\chapter{Теоријске основе}
\label{chp:teorijske-osnove}
\begingroup
\setlength{\emergencystretch}{3em}

\section{Критеријуми класификације}

Софтверска архитектура обухвата структуре система, односе њихових елемената и
својства значајна за разумевање и развој \cite{bass2012software}. Зато се не
одређује само на основу програмског оквира, инфраструктурног производа или
броја компонената. У овом раду варијанте се класификују према модулским,
процесним и распоредним границама, начину комуникације, власништву над подацима
и операционим последицама.

Посебно је важно разликовати три појма. \emph{Модулска граница} раздваја
одговорности и дозвољене зависности унутар кода. \emph{Граница процеса}
раздваја извршавање у различитим адресним просторима. \emph{Јединица
	распоређивања}\footnote{Јединица распоређивања је део система који се може
	изменити и поставити у рад без поновног постављања осталих делова. Она није
	нужно исто што и артефакт градње, процес или контејнер.} одређује шта се
објављује и поставља у рад као целина. Ове границе могу да се подударају, али
нису међусобно заменљиве.

\section{Монолит}

Монолит обједињује пословне компоненте у заједничкој јединици распоређивања:
измена једне компоненте типично захтева поновно распоређивање целине
\cite{lewis2014microservices,dragoni2017microservices}. Компоненте нису
независно распоредиве, али систем може имати јасну унутрашњу структуру и више
модула. Зато одсуство модула није услов ове класификације.

Број процеса није универзална дефиниција монолита. За класификацију је
пресудна заједничка јединица распоређивања и одсуство независно распоредивих
пословних компонената. Конкретна монолитна варијанта овог рада извршава се као
један апликациони процес по инстанци; то је својство имплементације које се
проверава у четвртом поглављу.

% Извор тврдње о конкретној варијанти: backend/media-monolith/docker-compose.yml
% и backend/media-monolith/Dockerfile.

\section{Модуларни монолит}

Модуларни монолит задржава заједничку јединицу распоређивања, али код организује
у слабо спрегнуте модуле са јасним одговорностима, јавним интерфејсима и
експлицитним зависностима \cite{su2024modularMonolith}. За његово разликовање
од полазног монолита потребни су докази о модулима, објављеним интерфејсима,
дозвољеним зависностима и очуваном заједничком распоређивању.

Аутоматска провера зависности није сама по себи дефиниција обрасца, већ
својство конкретне имплементације. Алат {\lat Spring Modulith} може из структуре
пакета извести модел модула, описати именоване интерфејсе и проверити циклусе,
приступ интерним пакетима и задата правила зависности
\cite{springModulith2026fundamentals}. Присуство алата није довољан доказ:
потребни су стварна правила и поновљив излаз њихове провере.

\section{Микросервисна архитектура}

Микросервисна архитектура организује систем као скуп кохезивних сервиса који
се извршавају у засебним процесима, комуницирају порукама и могу се независно
распоређивати \cite{dragoni2017microservices}. Величина сервиса није довољан
критеријум; одлучујуће су независне распоредне границе и уговори сарадње
\cite{lewis2014microservices}.

Комуникација преко границе процеса\footnote{Комуникација преко границе процеса
	одвија се посредством мрежног или другог међупроцесног механизма, а не
	непосредним позивом функције у истом адресном простору.} уводи могућност
недоступности учесника или везе и додатни трошак преноса
\cite{lewis2014microservices,dragoni2017microservices}. Зато су уговори,
поступање при отказу и могућност увида у рад система део поређења.
Децентрализовано управљање подацима јесте честа особина, али не захтева нужно
засебну физичку базу за сваки сервис \cite{lewis2014microservices}. У овом раду
се зато проверавају семантичко власништво над подацима, дозвољени приступи и
трансакционе границе.

\section{Синтеза критеријума}

Табела~\ref{tab:teorija-kriterijumi} сажима критеријуме који се касније
примењују на проверљиве артефакте сваке варијанте.

\begin{table}[!ht]
	\centering
	\small
	\begin{tabular}{|p{0.19\textwidth}|p{0.22\textwidth}|p{0.23\textwidth}|p{0.24\textwidth}|}
		\hline
		\textbf{Критеријум}    & \textbf{Монолит}                                                                                                          & \textbf{Модуларни монолит}                                                     & \textbf{Микросервиси}                                           \\
		\hline
		Јединица распоређивања & Заједничка целина; пословне компоненте нису независно распоредиве \cite{lewis2014microservices,dragoni2017microservices}. & Заједничка целина за све модуле \cite{su2024modularMonolith}.                  & Сервиси су независно распоредиви \cite{lewis2014microservices}. \\
		\hline
		Унутрашње границе      & Нису дефинициони услов; модули могу постојати.                                                                            & Јавни интерфејси и експлицитне зависности модула \cite{su2024modularMonolith}. & Уговори раздвајају сервисе \cite{lewis2014microservices}.       \\
		\hline
		Процесна граница       & У овом раду није дефинициони критеријум монолита.                                                                         & Заједнички апликациони процес \cite{su2024modularMonolith}.                    & Засебни процеси сервиса \cite{dragoni2017microservices}.        \\
		\hline
	\end{tabular}
	\caption{Критеријуми разграничења архитектуралних варијанти}
	\label{tab:teorija-kriterijumi}
\end{table}

За потребе миграције, \emph{граница по конвенцији} означава поделу изражену
пакетима, именовањем или документацијом без обавезне машинске провере.
\emph{Аутоматски проверљива граница} има формализовано правило и поновљив
поступак који открива његово нарушавање. Статичка провера архитектуралне
усаглашености пружа основу за такву контролу \cite{passos2010architectureConformance},
али способност алата постаје доказ тек када су правило, место извршавања и
исход провере документовани.

Архитектурални образац описује организацију и ограничења система, док алати и
инфраструктура омогућавају њихову реализацију или проверу
\cite{bass2012software}. Зато ни појединачан алат, број контејнера или избор
посредника за поруке сами не одређују архитектуру.
Класификација се заснива на стварним модулским, процесним и распоредним
границама и уговорима. На тој основи ИП1 пореди мерљиво понашање класификованих
варијанти, ИП2 њихове организационе и операционе последице, а ИП3 промене
између узастопних архитектуралних облика.
\endgroup

% ------------------------------------------------------------------------------
\chapter{Методологија}
\label{chp:metodologija}
\begingroup
\setlength{\emergencystretch}{3em}

Методологија операционализује три истраживачка питања. Поређење је
контролисано: систем под тестом, фактор, оптерећење, метрике и прикупљање
података одређују се пре тумачења исхода. Тако се ефекат варијанте не меша са
променом улаза, ресурса или окружења
\cite{jain1991art,kounev2020benchmarking}.

\section{Дизајн студије и систем под тестом}

Систем под тестом\footnote{Систем под тестом је скуп компонената чије се
	понашање пореди под контролисаним условима.} обухвата апликациони слој за
пријем и транскодовање видео-записа. Фактор је архитектурална варијанта са три
нивоа: монолит, модуларни монолит и микросервисна архитектура. За поређење се
користе извршиве варијанте са истим функционалним уговором и заједничким
протоколом.

% TODO: Пре званичног експеримента архивирати доказе да све три варијанте
% пролазе функционалне, инструменталне и протоколарне провере спремности.

Заједнички функционални ток обухвата отпремање видео записа, захтев за израду
варијанти, периодично читање статуса до коначног исхода и, када је сценаријом
предвиђено, преузимање резултата. Јавни уговор, пословна значења исхода и
ваљаност излазних видео записа проверавају се истим функционалним поступком за
све варијанте.

Примарна ресурсна граница је цео апликациони слој. У монолитној и
модуларно монолитној варијанти он обухвата њихову апликациону инстанцу, а у
микросервисној све унапред декларисане апликационе сервисе. Укупан буџет
процесора и меморије једног профила дели се унутар тог скупа; не додељује се у
целости сваком сервису. База података приказује се засебно као контекстуална
мера, док генератор оптерећења и мерна инфраструктура не улазе у примарни
обрачун.

% TODO: После верификације микросервисне имплементације додати један ауторски
% дијаграм који разјашњава границу система и ток мерења. Приказати генератор
% оптерећења, цео апликациони слој под заједничким ресурсним буџетом, базу као
% засебну контекстуалну меру, мерну инфраструктуру ван примарног обрачуна и
% ток захтев--обрада--исход--архива. Предлог натписа: „Граница система под
% тестом и ток прикупљања података“. Предлог ознаке:
% fig:metodologija-granica-i-tok. Коначна топологија сервиса мора потицати из
% проверене имплементације и конфигурације распоређивања.

\section{Променљиве, оптерећење и ресурсни профили}

Контролне променљиве су функционални ток и јавни уговор, тест-датотека,
тражени излази, сценарио оптерећења, верзија алата за транскодовање, хост и
конфигурација спољних зависности. Са фактором се намерно мењају модулске и
процесне границе, механизми комуникације и топологија распоређивања.

Свака варијанта испитује се у сваком ресурсном профилу. Профил је јединствен
третман који може истовремено одредити буџет процесора, меморије и број радних
нити; његов ефекат се зато не тумачи као изолован ефекат једне величине.

% TODO: После пилота унети и замрзнути коначне ресурсне профиле, број корисника,
% трајање увођења оптерећења, тест-датотеку, тражене излазе и временска
% ограничења. Не преносити подразумеване конфигурационе вредности као коначне.

Све варијанте користе исти хеш улазне датотеке, исте верзије зависности и
празно почетно стање. Не покрећу се паралелно. Потенцијални утицаји
температуре и фреквенције процесора, позадинских процеса, кеша, динамичког
превођења и редоследа варијанти умањују се стабилизацијом хоста, загревањем,
независним понављањима и уравнотеженим редоследом
\cite{georges2007statistically}.

\section{Операционе дефиниције метрика}

Метрика мора да одреди операцију, јединицу, мерни интервал и правило
укључивања \cite{jain1991art,kounev2020benchmarking}. Табела
\ref{tab:metodologija-metrike} даје дефиниције које се једнако примењују на
све варијанте.

\begin{table}[!ht]
	\centering
	\small
	\begin{tabular}{|p{0.20\textwidth}|p{0.45\textwidth}|p{0.20\textwidth}|}
		\hline
		\textbf{Мера}      & \textbf{Операциона дефиниција}                                                     & \textbf{Јединица}         \\
		\hline
		Пропусност         & Број потврђено завршених варијанти видео-записа подељен трајањем мерног интервала. & завршени излази у секунди \\
		\hline
		Време одзива       & Трајање захтева, засебно за сваку врсту захтева.                                   & милисекунде               \\
		\hline
		Латенција обраде   & Време коначног трајног исхода умањено за време креирања посла.                     & секунде по излазу         \\
		\hline
		Употреба процесора & Стопа утрошених процесорских секунди у граници апликационог слоја.                 & процесорска језгра        \\
		\hline
		Употреба меморије  & Радни скуп меморије у граници апликационог слоја.                                  & мегабајти                 \\
		\hline
	\end{tabular}
	\caption{Операционе дефиниције квантитативних метрика}
	\label{tab:metodologija-metrike}
\end{table}

Намеравани излаз је једна тражена варијанта видео-записа. Сваки намеравани
излаз у технички важећем покретању припада тачно једној од пет категорија:
потврђен \texttt{FINISHED}; потврђен \texttt{ERROR}; пословни неуспех пре
добијања потребног идентификатора; изостанак терминалног статуса до рока; или
технички изгубљен статус, када прикупљање не омогућава утврђивање исхода. Ове
категорије обухватају све намераване излазе; последња означава недоступност
мерења, а не успех или неуспех система. Пропусност користи само потврђене
\texttt{FINISHED} излазе, али се тумачи уз број и стопу свих категорија. Број
захтева у секунди није пословна пропусност.

Време одзива изводи се засебно по врсти захтева. Латенција обраде захтева
потврђене трајне временске ознаке креирања и коначног исхода; успешни и
неуспешни послови не спајају се у једну расподелу. Употреба процесора и
меморије изводи се из временских серија за цео апликациони слој и по његовим
компонентама. Узорци унутар покретања нису независна понављања.

% TODO: Пре мерења потврдити канонски извоз пословних исхода и трајних
% временских ознака, стабилну идентификацију компонената и тачан интервал
% узорковања; архивирати поступак израчунавања сваке метрике.

\section{Протокол извршавања и анализа}

Једно независно покретање, као јединица репликације, представља комбинацију
варијанте, ресурсног профила и редног броја понављања. Пре покретања база и
складиште се доводе у чисто стање, проверава се доступност компонената и
стабилизује хост. Потом се изводи загревање истог облика које не улази у
мерење. Загревање и независна извршавања су значајни за мерење система на
виртуелној машини {\lat Java} због динамичког превођења и варијабилности
извршавања \cite{georges2007statistically}.

Мерни интервал почиње непосредно пре увођења првог корисника и обухвата цео
сценарио и коначан период пражњења. Покретање се завршава када су сви
намеравани излази класификовани или када истекне унапред утврђен рок, али не
пре завршетка последњег планираног сценарија. Сви временски тренуци и коначни
исходи чувају се у архиви покретања.

% TODO: Пилотом одредити и пре званичне матрице замрзнути трајање загревања,
% коначни период пражњења, број независних понављања и поступак израчунавања
% интервала неизвесности.

Унутар сваког блока све три варијанте се извршавају на истом хосту, а њихов
редослед се уравнотежава тако да положај варијанте не буде стално исти.
Пилот-покретања се архивирају и искључују из примарне анализе. После пилота се
замрзавају протокол, параметри, распоред и поступак анализе.

Свако технички важеће и унапред планирано покретање остаје у анализи без
обзира на пословни исход. Недоступност појединачне метрике бележи се засебно и
не поништава остале проверљиве податке истог покретања. Недостајуће вредности
се не процењују.

Метрике се најпре сажимају на нивоу покретања, а затим се приказују независна
понављања и њихова варијабилност \cite{jain1991art,kounev2020benchmarking,
	georges2007statistically}. Примарни пар чине покретања истог ресурсног профила
и истог унапред пријављеног блока понављања. Монолит је референца, па су
примарне процене апсолутна разлика „модуларни монолит минус монолит” и однос
„модуларни монолит подељен монолитом”, односно исте две процене за
микросервисе према монолиту. Процене се изводе на нивоу покретања. Покретања
која нису изведена као унапред пријављен пар не третирају се накнадно као
упарена. Профили се не спајају у јединствен ранг.

\section{Квалитативно поређење и евиденција миграције}

За ИП2 се на све варијанте примењују исти унапред одређени поступци:
\begin{itemize}
	\item организација кода описује се мапом одговорности, јавних интерфејса,
	      зависности и машински проверљивих правила;
	\item распоређивање се пореди за исту унапред одређену класу измене, уз
	      бележење исхода поступка;
	\item понашање при отказу испитује се истом унапред одређеном пробом, уз
	      исти прозор посматрања и критеријум опоравка.
\end{itemize}

За поређење распоређивања унапред се утврђују полазно стање, класа измене,
обухваћене јединице, предуслови и кораци поступка. За пробе отказа утврђују се
почетни подаци, уочљиви окидач, трајање прекида, прозор посматрања, критеријум
опоравка и број понављања. Ове вредности једнаке су за све варијанте и чувају
се са исходима.

Пробе обухватају отказ апликационе компоненте током активне обраде, отказ базе
током прихватања и обраде и неважећи или недоступан улаз. За сваку пробу
бележе се посматрани исходи у унапред одређеном прозору и испуњеност
критеријума опоравка.

% TODO: После имплементације микросервиса прецизирати симетричан избор сервиса
% за пробу отказа без мењања пословног предуслова, трајања и критеријума
% опоравка; потом архивирати поступак и исход сваке пробе.

За ИП3 је јединица анализе \emph{документована миграциона промена}: једна
семантичка промена са јединственим идентификатором, циљем, полазном и циљном
варијантом, категоријом и доказним ревизијама или путањама. Више датотека које
остварују исти циљ групишу се у једну промену, а премештање без промене значења
бележи се засебно. Дневник се води за обе етапе миграције по истом правилу и
садржи доказе довољне за проверу врсте и обима промене. Број датотека и линија
може бити описни атрибут, али није мера труда. Описни обим се унапред одређује
као број документованих јединица по категорији и скуп категорија захваћених
артефаката; не представља време, трошак или уложени рад.

% TODO: Довршити и архивирати дневнике обе миграције са јединственим
% идентификаторима и правилима груписања; без временске евиденције не изводити
% човек-сате, календарско трајање или трошак.

\section{Синтеза, валидност и поновљивост}

Препоруке се изводе повезивањем налаза ИП1--ИП3 са практичним приоритетом,
испитаним профилом, ограничењем и условом примене. Не формира се јединствена
оцена без образложених тежина. Препорука је оправдана само ако је њен
предуслов непосредно мерен или посматран.

Конструктну валидност угрожава замена пословне пропусности бројем захтева или
временске употребе ресурса тренутним снимком; зато су мере операционо раздвојене.
Унутрашњу валидност угрожавају редослед, променљиво стање хоста и истовремена
промена више својстава ресурсног профила. Уравнотежен редослед, чисто стање,
загревање, независна понављања и тумачење профила као целине умањују ове
утицаје \cite{georges2007statistically}.

Спољашња валидност ограничена је доменом транскодовања, изабраним
технолошким скупом, синтетичким улазом, дефинисаним оптерећењем и локалним
окружењем. Налази се односе на конкретне имплементације под испитаним условима.
Закључну валидност ограничавају мали број независних понављања, варијабилност
међу покретањима, више исхода и поређења и још неодређен поступак израчунавања
интервала неизвесности. Зато се приказују сва понављања и неизвесност, а
одсуство уочене разлике не тумачи се као доказ једнакости.

Поновљивост захтева архиву која за свако покретање повезује варијанту,
ревизију имплементације, ресурсни профил, улаз, верзију протокола, сирове
излазе, коначне пословне исходе и поступак анализе.

% TODO: Пре писања резултата проверити потпуност архиве и следљивост сваке
% табеле и графикона до независних покретања и сирових података.

Наредна три поглавља примењују утврђене критеријуме и метод на монолитну,
модуларно-монолитну и микросервисну варијанту и на промене између њих.
\endgroup

% ------------------------------------------------------------------------------
\chapter{Монолитна архитектура}
\label{chp:monolitna-arhitektura}
% TODO: Написати из проверене грађе у docs/thesis/thesis-notes/04-monolith.md.

% TODO: Извор дијаграма је проверена имплементација и docs/app-specs/monolith-spec.md.
\begin{figure}[!ht]
	\draftfigureplaceholder{КОМПОНЕНТНИ И КОНТЕЈНЕРСКИ ДИЈАГРАМ МОНОЛИТА}{У
	граници једне {\lat Spring Boot} апликације приказати компоненте
	\texttt{video} и \texttt{notification}, њихове портове и смер комуникације.
	Приказати клијентски {\lat REST}/{\lat SSE} приступ, једну базу
	{\lat PostgreSQL}, локално складиште фајлова и спољне процесе
	{\lat FFprobe}/{\lat FFmpeg}. Разликовати диспеч послова преко
	{\lat PostgreSQL LISTEN/NOTIFY} од исходних догађаја који се преносе
	{\lat Spring ApplicationEvents} механизмом ка компоненти за обавештења;
	означити ред послова, скуп радних нити, упис статуса и слање обавештења
	клијенту.}
	\caption{Компоненте монолитне варијанте и њене спољне зависности}
	\label{fig:monolit-komponente-zavisnosti}
\end{figure}

% ------------------------------------------------------------------------------
\chapter{Модуларни монолит}
\label{chp:modularni-monolit}
% TODO: Написати након архивирања проверљивих артефаката миграције.

% TODO: Регенерисати components.puml провером ModularityTests и извести га под thesis/figures/.
\begin{figure}[!ht]
	\draftfigureplaceholder{ГРАФ ЗАВИСНОСТИ МОДУЛА}{Унети граф који алат
	{\lat Spring Modulith} генерише из проверене структуре. Приказати модуле
	\texttt{config}, \texttt{shared}, \texttt{video} и
	\texttt{notification}, као и дозвољене ивице
	\texttt{video}~$\rightarrow$~\texttt{shared},
	\texttt{video}~$\rightarrow$~\texttt{config} и
	\texttt{notification}~$\rightarrow$~\texttt{shared}. Посебно означити да
	\texttt{notification} слуша објављене догађаје модула \texttt{video},
	уместо да зависи од његових интерних типова, и издвојити именовани
	интерфејс \texttt{video.published}.}
	\caption{Граф модула у модуларном монолиту генерисан алатом {\lat Spring Modulith}}
	\label{fig:modularni-monolit-graf-modula}
\end{figure}

% ------------------------------------------------------------------------------
\chapter{Микросервисна архитектура}
\label{chp:mikroservisna-arhitektura}
% TODO: Поглавље написати после имплементације микросервисне варијанте и
% архивирања артефаката друге етапе миграције; не представљати пројектована
% својства као потврђене резултате.

% TODO: Попунити искључиво према коначно имплементираној и провереној топологији.
\begin{figure}[!ht]
	\draftfigureplaceholder{КОМПОНЕНТНИ И КОНТЕЈНЕРСКИ ДИЈАГРАМ МИКРОСЕРВИСА}{После
	имплементације приказати стварне границе независно распоредивих сервиса,
	власништво над ендпоинтима и подацима и све инфраструктурне контејнере.
	Дијаграм треба да обухвати улазну тачку јавног {\lat API}-ја, сервис који
	управља видео-садржајем и транскодовањем, сервис за обавештења, њихове
	базе или шеме, изабрани брокер порука и објектно складиште, као и циљеве
	прикупљања метрика. Стрелицама разликовати синхроне захтеве, асинхроне
	поруке и приступ фајловима. Не уносити називе, број сервиса ни односе који
	нису потврђени коначним кодом и конфигурацијом распоређивања.}
	\caption{Границе сервиса и инфраструктурне зависности микросервисне варијанте}
	\label{fig:mikroservisi-granice-infrastruktura}
\end{figure}

% TODO: Секвенцу ускладити са коначним уговорима порука и јавним API-јем.
\begin{figure}[!ht]
	\draftfigureplaceholder{СЕКВЕНЦИЈСКИ ДИЈАГРАМ ДИСТРИБУИРАНЕ ОБРАДЕ}{После
	имплементације приказати учеснике који стварно постоје и следећи логички
	ток: клијент отпрема видео и тражи израду варијанте; надлежни сервис чува
	објекат и метаподатке, бележи посао и асинхроно га објављује; сервис за
	обраду преузима поруку и фајл из објектног складишта, покреће
	{\lat FFmpeg}, чува резултат и објављује догађај исхода; сервис за
	обавештења бележи исход и, ако тај механизам остане у коначној варијанти,
	шаље {\lat SSE} обавештење. Приказати и клијентову проверу статуса, границе
	трансакција и гране успеха и отказа, без унапред задате гаранције испоруке.}
	\caption{Секвенца асинхроне израде варијанте и обавештавања у микросервисној архитектури}
	\label{fig:mikroservisi-sekvenca-obrade}
\end{figure}

% ------------------------------------------------------------------------------
\chapter{Резултати и дискусија}
\label{chp:rezultati}
% TODO: Поглавље написати тек после званичне матрице и архивирања примарних
% артефаката, упита и поступка агрегације.

% TODO: Тип графикона и статистички сажетак изабрати тек после прегледа архивираних података.
\begin{figure}[!ht]
	\draftfigureplaceholder{УПОРЕДНЕ ГРАФИКОНЕ РЕЗУЛТАТА}{На основу званично
		архивираних покретања припремити вишепанелни приказ за три архитектуралне
		варијанте и четири сложена профила ресурса. За сваку комбинацију приказати
		сва три понављања као појединачне тачке и јасно представити
		варијабилност; не приказивати само средњу вредност. Раздвојити панеле за
		пропусност, време одзива, латенцију тока обраде, искоришћеност процесора и
		потрошњу меморије, са јединицама и истом легендом. Избор распона оса,
		мере средишње вредности и интервала или другог показатеља варијабилности
		засновати на коначној расподели података, без унапред унетих вредности.}
	\caption{Поређење архитектуралних варијанти по профилу ресурса, понављању и метрици}
	\label{fig:rezultati-varijabilnost-profili}
\end{figure}

% ------------------------------------------------------------------------------
\chapter{Преглед литературе}
\label{chp:pregled-literature}
% TODO: Написати искључиво из проверених извора и позиционирати стварни допринос рада.

% ------------------------------------------------------------------------------
\chapter{Закључак}
\label{chp:zakljucak}
% TODO: Написати после резултата; не уносити коначне налазе или препоруке унапред.

% ------------------------------------------------------------------------------
% Literatura
% ------------------------------------------------------------------------------
\literatura

% ==============================================================================
% Završni deo teze i prilozi
\backmatter
% ==============================================================================

% TODO: Додати биографију кандидата тек када аутор достави проверене податке.

\end{document}
